% !PRINTED 749
\weberpage{749}

% [sic] -- the source uses $\delta_1$ and $\delta_2$ both for scalar lengths and for intervals; the printed symbols are preserved.
Wir bezeichnen mit $\alpha_1$, $\beta_1$ die beiden \gesperrt{fr\"uhesten} Zahlen der Reihe $\Omega$, die in dem Intervalle $\delta$ liegen, nehmen $\alpha_1 < \beta_1$ an, und setzen $\delta_1 = \beta_1 - \alpha_1$, so dass $\delta_1 = (\alpha_1, \beta_1)$ ein Theil des Intervalles $\delta$ ist.

Nun bezeichnen wir ebenso mit $\alpha_2$, $\beta_2$ die beiden fr\"uhesten Zahlen von $\Omega$, die \gesperrt{im Inneren} des Intervalles $\delta_1$ (mit Ausschluss der Grenzen) liegen, setzen $\alpha_2 < \beta_2$ voraus, und setzen $\delta_2 = \beta_2 - \alpha_2$.

Auf diese Weise k\"onnen wir fortfahren und erhalten eine unbegrenzte Reihe von Intervallen:
\[
\delta,\ \delta_1,\ \delta_2,\ \delta_3,\ \dots,
\]
deren jedes alle folgenden einschliesst, und zwei Reihen von Zahlen:
\[
\begin{gathered}
\alpha,\ \alpha_1,\ \alpha_2,\ \dots \\
\beta,\ \beta_1,\ \beta_2,\ \dots
\end{gathered}
\]
die, mit etwaiger Ausnahme der ersten, $\alpha$ und $\beta$, alle der Reihe $\Omega$ angeh\"oren. Die $\alpha, \alpha_1, \alpha_2, \dots$ bilden eine wachsende, die $\beta, \beta_1, \beta_2, \dots$ eine fallende Zahlenreihe, und zugleich ist jedes $\alpha$ kleiner als jedes $\beta$.

1. \gesperrt{Daraus ergiebt sich, dass die Zahlen $\alpha_\nu$ eine obere Grenze $a$, die Zahlen $\beta_\nu$ eine untere Grenze $b$ haben, und dass $a$ jedenfalls nicht gr\"osser als $b$ ist.} (Vgl. Bd. I, \S.\,38, 1.)

Es kann aber m\"oglicherweise $a = b$ sein.

Aus der Bildungsweise der Intervalle $\delta, \delta_1, \delta_2, \dots$ geht noch Folgendes hervor:

2. \gesperrt{Wenn irgend eine Zahl $\omega$ der Reihe $\Omega$ in dem Intervalle $\delta_\nu$, mit Ausschluss seiner Grenzen $\alpha_\nu, \beta_\nu$, liegt, so ist $\omega$ in der Reihe $\Omega$ sp\"ater als das derselben Reihe angeh\"orige Zahlenpaar $\alpha_\nu, \beta_\nu$.}

Denn $\alpha_\nu, \beta_\nu$ waren ja die beiden \gesperrt{fr\"uhesten} im Intervalle $\delta_{\nu-1}$ gelegenen Zahlen von $\Omega$. Daraus folgt:

3. \gesperrt{Die der Reihe $\Omega$ angeh\"origen Zahlenpaare $\alpha_\nu, \beta_\nu$ sind um so sp\"atere Glieder der Reihe $\Omega$, je gr\"osser der Index $\nu$ ist, und da wir angenommen haben, die Reihe der Intervalle $\delta_\nu$ breche nicht ab, so k\"onnen wir $\alpha_\nu, \beta_\nu$ f\"ur ein hinl\"anglich grosses $\nu$ beliebig weit in der Reihe $\Omega$ hinausr\"ucken lassen.}
